\documentclass[letterpaper,11pt]{moderncv}
\moderncvstyle{banking}
\moderncvcolor{black}

\usepackage[letterpaper, margin=0.8in]{geometry}
\usepackage{xcolor}
\usepackage{fontspec}
\setmainfont{Times New Roman}
\usepackage[unicode]{hyperref}

\name{Saeyoung}{Kim}
\address{Cambridge, MA, USA}{}{}
\phone[mobile]{+1 (351) 322 5510}
\email{saeyoung\_kim@uml.edu}
\social[linkedin]{saeyoung-macx-kim}

\recipient{Hiring Team}{ASML\\Austin, TX}
\date{\today}
\opening{Dear Hiring Team,}
\closing{Sincerely,}

\begin{document}

\makelettertitle

I'm writing about the Upgrade Install and Relocation Engineer position (J-00336327) in Austin. Most of my recent work has been assembling, testing, and qualifying hardware in cleanroom environments, so this role is a natural fit.

At Hanwha Systems, I spent over a year on military SAR satellite programs running AIT campaigns on Flight Model hardware. Day to day, that meant setting up EGSE test configurations, writing and executing integration test procedures, chasing down anomalies when things didn't behave as expected, and documenting everything along the way. The work was procedural and deadline-driven, and I got used to coordinating across different engineering teams and with customers to keep things moving.

Before that, my PhD was very hands-on. I designed circuits, laid out PCBs, did cleanroom fabrication (photolithography, electroplating, the whole process), and then integrated and tested the finished sensor systems. When something broke, I had to figure out whether it was a hardware problem, a firmware bug, or a process issue. That kind of troubleshooting is second nature to me now.

I'd like to bring that experience to ASML's lithography install work. Happy to talk more if it seems like a good fit. Thanks for your time.

\makeletterclosing

\end{document}
